\documentclass[11pt,a4paper]{article}

\usepackage[utf8]{inputenc}
\usepackage{fontspec}
\setmainfont{Times New Roman}
\usepackage[main=russian, english]{babel} % Подключение русского языка
\usepackage[left=1.5cm, right=1.5cm, top=1.5cm, bottom=1.5cm]{geometry}
\usepackage[hidelinks]{hyperref}
\usepackage{xcolor}
\hypersetup{colorlinks=true, linkcolor=black, filecolor=black, urlcolor=blue!70!black}
\usepackage{enumitem}
\usepackage{titlesec}

\titleformat{\section}{\large\bfseries\uppercase}{}{0em}{}[\titlerule\vspace{0.5ex}]
\titlespacing{\section}{0pt}{10pt}{5pt}
\setlength{\parindent}{0pt}

\begin{document}

% -------------------- HEADER --------------------
\begin{center}
    {\Huge \textbf{Игорь Полуянов}} \\[0.3cm]
    {\huge Backend Go-разработчик} \\[0.3cm]
    Санкт-Петербург | +7 992 200 8123 | \href{mailto:poluyanovia@yandex.ru}{poluyanovia@yandex.ru} \\
    \href{https://github.com/PoluyanbIch}{github.com/PoluyanbIch} | Telegram: \href{https://t.me/PoluyanbIch}{@PoluyanbIch}
\end{center}

% -------------------- SUMMARY --------------------
\section{Обо мне}
\textbf{Backend-разработчик} с опытом проектирования распределенных систем на \textbf{Go}, \textbf{Python} и \textbf{Java}. Имею практические навыки разработки микросервисной архитектуры и REST/gRPC API. Фокусируюсь на качестве кода (покрытие тестами 80\%+), оптимизации баз данных и эффективном использовании многопоточности в Go. Быстро погружаюсь в новые технологии и бизнес-логику проектов.

% -------------------- EDUCATION --------------------
\section{Образование}
\textbf{Университет ИТМО} \hfill \textit{Ожидаемый год окончания: 2028} \\
Бакалавриат -- Факультет программной инженерии и компьютерной техники (ПИиКТ) \\
\textit{Специальность: Программная инженерия, Системное и прикладное программное обеспечение}

% -------------------- SKILLS --------------------
\section{Навыки}
\textbf{Языки программирования:} Golang, Python, Java, SQL, Bash \\
\textbf{Архитектура и API:} gRPC, REST API, Микросервисы, Гексагональная архитектура, SOLID, ООП \\
\textbf{Базы данных и брокеры:} PostgreSQL, Redis, MongoDB, NATS, Kafka \\
\textbf{Инструменты и DevOps:} Git, Docker, docker-compose, CI/CD, Swagger/OpenAPI \\
\textbf{Тестирование и мониторинг:} Unit-тестирование (Testify, go test), Grafana, Prometheus, Kibana

% -------------------- WORK EXPERIENCE --------------------
\section{Опыт работы}
\textbf{RIO SOFT} -- \textit{Пермь, Россия} \hfill \textit{Май 2025 -- Сен 2025} \\
\textit{Стажер Backend-разработчик}
\begin{itemize}[leftmargin=0.5cm, itemsep=2pt, topsep=2pt]
    \item Разработал и внедрил 3 новые ключевые функции для внутреннего Telegram-бота с использованием \textbf{Python} и \textbf{REST API}, автоматизировав ежедневную отчетность для команды из 50+ сотрудников.
    \item Занимался поддержкой и рефакторингом legacy веб-приложения на \textbf{Java} и \textbf{Spring MVC}: исправлял баги, оптимизировал логику работы методов, писал unit-тесты.
    \item Оптимизировал SQL-запросы к \textbf{PostgreSQL} в нагруженных частях приложения, за счет чего время выполнения критических выборок снизилось примерно на 40\%.
    \item Унифицировал подход к передаче данных, внедрив DTO вместо прямой работы со слоем сущностей (Entity), что снизило нагрузку на сеть и ускорило ответы API до 40\%.
    \item Работал по методологии Agile/Scrum, участвовал в код-ревью и версионировании кодовой базы с помощью Git.
\end{itemize}

% -------------------- PROJECTS & HACKATHONS --------------------
\section{Проекты и Хакатоны}

\textbf{Микросервисная поисковая система для xkcd.com} | \textit{Pet-проект} \hfill \textit{2025}
\begin{itemize}[leftmargin=0.5cm, itemsep=0pt, topsep=2pt]
    \item Спроектировал бэкенд на \textbf{Go} (гексагональная архитектура), состоящий из 5 микросервисов (API Gateway, Search, AAA и др.) с синхронным взаимодействием по \textbf{gRPC} и асинхронным через брокер \textbf{NATS}.
    \item Разработал масштабируемый полнотекстовый поиск на базе кастомного инвертированного индекса в \textbf{PostgreSQL} с его автоматической перестройкой в фоновом режиме.
    \item Внедрил \textbf{JWT-аутентификацию} (Access/Refresh), написал кастомные middleware для защиты эндпоинтов (rate/concurrency limiter), полностью покрыл бизнес-логику модульными тестами (\textbf{Testify}) и контейнеризировал проект через \textbf{Docker Compose}.
\end{itemize}

\vspace{4pt}
\textbf{Сервис отслеживания настроений в реальном времени} | \textit{Хакатон Мэра Москвы «ЛЦТ 2025»} \hfill \textit{2025}
\begin{itemize}[leftmargin=0.5cm, itemsep=0pt, topsep=2pt]
    \item Спроектировал бэкенд для аналитики на \textbf{Golang}, применив конкурентные паттерны Go (goroutines, channels) для параллельной обработки входящего потока симулированных клиентских событий.
\end{itemize}

\vspace{4pt}
\textbf{Чат-бот для мессенджера} | \textit{Хакатон MAX \& VK Education} \hfill \textit{2025}
\begin{itemize}[leftmargin=0.5cm, itemsep=0pt, topsep=2pt]
    \item Написал диалогового чат-бота для мессенджера MAX на \textbf{Python}, обрабатывающего 500+ пользовательских взаимодействий ежедневно.
\end{itemize}

\vspace{4pt}
\textbf{Алгоритмический робот для лабиринта} | \textit{MTS True Tech Champ 2025} \hfill \textit{2025}
\begin{itemize}[leftmargin=0.5cm, itemsep=0pt, topsep=2pt]
    \item Реализовал алгоритмы поиска пути на графах на \textbf{Python}, позволившие роботу пройти лабиринт 50x50 менее чем за 30 секунд без столкновений.
\end{itemize}

\vspace{4pt}
\textbf{Прочие хакатоны и достижения}
\begin{itemize}[leftmargin=0.5cm, itemsep=0pt, topsep=2pt]
    \item Альфа Кейс-чемпионат, Кейс-чемпионат MIR, Changellenge >> CUP RUSSIA, Changellenge >> CUP IT.
    \item Выполнил исследовательский проект по очистке и анализу данных (\textbf{Python}) для датасетов объемом 10 000+ записей (DANO \& MTS, 2023).
\end{itemize}

% -------------------- EDUCATIONAL COURSES --------------------
\section{Курсы и сертификаты}
\begin{itemize}[leftmargin=0.5cm, itemsep=2pt, topsep=2pt]
    \item \textbf{Т-Образование:} Алгоритмы и структуры данных \hfill \textit{2026}
    \item \textbf{Яндекс:} Тренировки по алгоритмам \hfill \textit{2026}
    \item \textbf{Yadro Education:} Разработка микросервисных приложений на Go \hfill \textit{2025}
    \item \textbf{Яндекс Лицей:} Основы промышленного программирования на Python \hfill \textit{2023 -- 2024}
    \item \textbf{Яндекс Лицей:} Основы программирования на Python \hfill \textit{2022 -- 2023}
\end{itemize}

\end{document}